% !TEX TS-program = pdflatex
% David - Metsäalan harjoittelija CV (suomi)

\documentclass[a4paper,11pt]{article}

\IfFileExists{cv-style.sty}{\usepackage{cv-style}}{\usepackage{../../shared/cv-style}}

\cvlang{fi}
\cvrole{Metsäalan harjoittelija}

\begin{document}
\small
\setlength{\cvlabelw}{3.9cm}
\titlespacing{\section}{0pt}{0.55em}{2pt}

\begin{cvheader}
\par{\centering
		{\Huge \textsc{David Adegbola}}\\[0.4em]
		{\Large Metsäalan harjoittelija}
	\bigskip\par}

\hrule
\vspace{0.5em}
\cvline[2.4cm]{\cvt{location}:}{Joensuu, Suomi (voin työskennellä missä päin Suomea tahansa)}
\cvline[2.4cm]{\cvt{phone}:}{+358 46 800 7210}
\cvline[2.4cm]{\cvt{email}:}{\href{mailto:adavidadegbola@gmail.com}{adavidadegbola@gmail.com}}
\end{cvheader}

\section{\cvt{summary}}
\begin{tabular}{p{0.9\textwidth}}
 Hei! Olen David, 24-vuotias kestävän metsäbiotalouden (SuFoBE) opiskelija Itä-Suomen yliopistossa Joensuussa. Näin metsäteollisuuden ensimmäisen kerran sisältä päin UPM Kaukaan sellutehtaalla, kun olin asentajana sähkösuodattimen uudelleenrakennushankkeessa. Työnantaja arvioi työni kiitettäväksi ja kuvasi minua halukkaaksi oppimaan ja ottamaan vastuuta. Sen työn takia valitsin metsäbiotalouden.\\\\

 Taustani on tietotekniikassa, ja Vaasa Startup Factoryssa selvitin käyttäjien tarpeita haastatteluilla ja soittamalla yrityksille. Viihdyn raskaassa fyysisessä ulkotyössä. Haen metsäalan harjoittelu-, kesä- ja maastotöitä. Voin tehdä kokoaikatyötä, aloittaa heti ja matkustaa. Suomen kielen taitoni on vielä perustasolla, ja opiskelen sitä aktiivisesti.
\end{tabular}

\section{\cvt{education}}
\cvline{2026--nykyhetki:}{B.Sc, kestävä metsäbiotalous, Itä-Suomen yliopisto (UEF), Joensuu}
\cvline{2024--2026:}{B.Eng-opinnot, tietotekniikka, Vaasan ammattikorkeakoulu (VAMK), kesken}
\cvline{2019--2022:}{OND, tietojenkäsittely, The Polytechnic, Ibadan, Nigeria}

\section{\cvt{skills}}
\cvline{Selvitystyö:}{Käyttäjätarpeiden selvitys haastatteluilla ja yrityksille soittamalla}
\cvline{Työmaa:}{Työturvallisuusohjeiden noudattaminen; kulmahiomakone ja nostolaitteet}
\cvline{Työtavat:}{Tiimityö; vastuullisuus; raskas fyysinen ulkotyö}
\cvline{Haluan oppia:}{Paikkatieto (GIS), metsän mittaus, raivaussaha, moottorisaha}
\cvline{\cvt{cards}:}{Työturvallisuuskortti}
\cvline{\cvt{languages}:}{Suomi (perustaso, opiskelen aktiivisesti); englanti (sujuva); yoruba (äidinkieli)}

\section{\cvt{experience}}
\begin{cventry}{elo 2025--loka 2025}{\job{Asentaja}{Axioma Servisas UAB}}
  UPM Kaukaan sellutehdas, Lappeenranta\par
  Asennustyöt sähkösuodattimen uudelleenrakennushankkeessa metsäteollisuuden tehtaalla.\par
  Kulmahiomakoneen käyttö ja nostolaitteiden käyttö.\par
  Noudatin tehtaan tiukkoja turvallisuusohjeita, pidin työmaan siistinä ja seurasin työkaluja ja materiaaleja.\par
  Työnantajan arvio: kiitettävä. Vastuuntuntoinen ja tunnollinen työntekijä, halukas oppimaan ja ottamaan vastuuta.
  \sotagrow{metsäteollisuus,sellutehdas,työturvallisuus,nostolaitteet,oppimishalu}
\end{cventry}
\begin{cventry}{kesä 2026--elo 2026}{\textbf{Founder in Residence}, Vaasa Startup Factory, Vaasa: selvitin käyttäjätarpeita haastatteluilla ja yrityksille soittamalla; rakensimme ja esittelimme tiimin kanssa hetkin.appin.}
\end{cventry}
\begin{cventry}{tammi 2026--nykyhetki}{\textbf{Tekoälykehittäjä (omat projektit)}}
  Rakennan käytännön työkaluja sekä kehitän ja automatisoin sovelluksia tekoälyn avulla (Claude Code): kirjoitan määrittelyn, ohjaan agenttia ja tarkistan tuloksen.\par
  Kirjoitin sääntötiedoston, joka pitää agentin luotettavana: kiinteät faktat, kielletyt väitteet ja tarkistukset, jotka se tekee ennen jokaista tehtävää.\par
  Rakennan tekoälyagenttien avulla metsäaiheista oppimispeliä lapsille (kesken).\par
  \faGithub{} \href{https://github.com/david-adegbola}{github.com/david-adegbola}
  \sotagrow{ai-automation,claude-code,agent-rules,github-actions,automation}
\end{cventry}
\begin{cventry}{huhti 2025--elo 2025}{\textbf{Lehdenjakaja}, Finnbox ja SSM, Vaasa: aamureitit yksin ulkona säällä kuin säällä, aina ajallaan.}
\end{cventry}
\begin{cventry}{marras 2024--elo 2026}{\textbf{AV-vapaaehtoinen}, Vaasa City Church, Vaasa: ääni- ja videolaitteet joka viikonloppu.}
\end{cventry}
\begin{cventry}{syys 2016--syys 2020}{Myynti- ja asiakaspalvelutyötä Nigeriassa (Desire Feeds, Bunmi Pastries).}
\end{cventry}

\section{\cvt{hobbies}}
Kuntosali, asioiden korjaaminen ja valokuvaus.

\section{\cvt{references}}
\cvt{refsavailable}

\end{document}
